% Source: https://artofproblemsolving.com/community/c3h1910919_lines_of_symmetry
% Stable identity: aops-topic-1910919-q01
$\boxed{n}$. A line of symmetry of a regular $n$-gon passes through its centre, and a reflection maps vertices to vertices, so it is determined by where it sends a fixed vertex: it can send it to any of the $n$ vertices, giving at most $n$ reflections. Conversely, each line through the centre and a vertex, and each line through the centre and the midpoint of a side, is a line of symmetry.
\begin{asy}
usepackage("amsmath");
size(3.5cm);
defaultpen(fontsize(10pt));
int n=6;pair[] V;for(int k=0;k<n;++k)V.push(dir(90+360*k/n));draw(V[0]--V[1]--V[2]--V[3]--V[4]--V[5]--cycle);
for(int k=0;k<3;++k){draw(1.2*V[k]--1.2*V[k+3],dashed);pair m=(V[k]+V[k+1])/2;draw(1.3*m--(-1.3)*m,dotted);}
\end{asy}

For $n$ odd these are the $n$ lines through a vertex and the midpoint of the opposite side; for $n$ even there are $\dfrac n2$ through opposite vertices (dashed) and $\dfrac n2$ through midpoints of opposite sides (dotted). Either way there are exactly $n$.
